
\begin{tikzpicture}[
    >=Stealth,
    sensor/.style={
        circle,
        draw=black!85,
        fill=gray!25,
        line width=1pt,
        minimum size=0.58cm
    },
    source/.style={
        circle,
        draw=black!85,
        fill=orange!75,
        line width=1.1pt,
        minimum size=0.58cm
    },
    axis/.style={
        ->,
        line width=0.75pt,
        draw=black!60
    },
    baseline/.style={
        <->,
        draw=black!70,
        line width=0.7pt
    },
    pathRone/.style={
        ->,
        line width=1.2pt,
        draw=UdeAGreen!85!black
    },
    pathRtwo/.style={
        ->,
        line width=1.2pt,
        draw=blue!80!black
    },
    reference/.style={
        dashed,
        line width=0.9pt,
        draw=black!65
    },
    projection/.style={
        dotted,
        line width=0.65pt,
        draw=black!50
    },
    wavefront/.style={
        line width=0.95pt,
        draw=orange!90!black
    }
]

% ==============================================================================
% 1. RECUADRO MATEMÁTICO
% ==============================================================================

\node[
    draw=black!70,
    line width=0.8pt,
    fill=white,
    rounded corners=4pt,
    inner sep=9pt,
    anchor=west
] at (-6.1, 3.1) {
    \small
    \begin{tabular}{@{} l @{}}
        \textbf{Geometría angular y PDoA:} \\[5pt]

        Aproximación de campo lejano. \\[5pt]

        Diferencia de camino:
        \ $\Delta r = r_1-r_2=d\sin\theta$ \\[5pt]

        Diferencia de fase:
        \ $\Delta\Phi=
        \dfrac{2\pi f_0}{c(T)}\Delta r$ \\[7pt]

        \textcolor{UdeAGreen!85!black}{
            \textbf{Condición de no-aliasing espacial:}
        } \\[3pt]

        $d\leq\dfrac{\lambda_0}{2}
        =\dfrac{c(T)}{2f_0}$
    \end{tabular}
};

% ==============================================================================
% 2. ARREGLO LINEAL DE RECEPTORES
% ==============================================================================

% Centro geométrico del arreglo.

\coordinate (Orig) at (3.8, 0);

% Eje transversal.
% Se acorta su extremo derecho para mantener la etiqueta
% dentro de los márgenes de la figura.

\draw[axis]
    (0.35, 0) -- (7.2, 0)
    node[right=3pt, font=\small] {$x$ (transversal)};

% Eje normal al arreglo.

\draw[axis]
    (3.8, -0.65) -- (3.8, 7.25)
    node[above=3pt, font=\small] {$y$ (normal al arreglo)};

% Posiciones de los receptores.
% La separación se exagera deliberadamente para facilitar
% la lectura de la geometría.

\coordinate (R1) at (1.7, 0);
\coordinate (R2) at (5.9, 0);

\node[sensor] (NodeR1) at (R1) {};

\node[
    below=12pt,
    font=\small
] at (R1) {
    $R_1\left(-\frac{d}{2},0\right)$
};

\node[sensor] (NodeR2) at (R2) {};

\node[
    below=12pt,
    font=\small
] at (R2) {
    $R_2\left(+\frac{d}{2},0\right)$
};

% ==============================================================================
% 3. COTA DE SEPARACIÓN ENTRE RECEPTORES
% ==============================================================================

% Se baja la cota para evitar cualquier solapamiento con
% las etiquetas de los receptores.

\draw[baseline]
    (1.7, -1.65) --
    node[below=3pt, font=\small]
    {$d=\|\vec{r}_2-\vec{r}_1\|$}
    (5.9, -1.65);

% Líneas auxiliares de la cota.

\draw[projection]
    (1.7, 0) -- (1.7, -1.65);

\draw[projection]
    (5.9, 0) -- (5.9, -1.65);

% ==============================================================================
% 4. FUENTE ACÚSTICA
% ==============================================================================

% Fuente puntual situada por encima y a la derecha del arreglo.

\coordinate (S) at (6.6, 6.35);

\node[source] (NodeS) at (S) {};

\node[
    above right=4pt,
    font=\small\bfseries
] at (NodeS) {Fuente $S(x,y)$};

% ==============================================================================
% 5. FRENTES DE ONDA
% ==============================================================================

% Arcos centrados en la fuente.
% Representan la propagación desde la fuente hacia los receptores.

\draw[wavefront]
    ($(S)+(205:0.48)$)
    arc[start angle=205, end angle=265, radius=0.48];

\draw[wavefront]
    ($(S)+(205:0.78)$)
    arc[start angle=205, end angle=265, radius=0.78];

\draw[wavefront]
    ($(S)+(205:1.08)$)
    arc[start angle=205, end angle=265, radius=1.08];

% ==============================================================================
% 6. TRAYECTORIAS ENTRE LA FUENTE Y LOS RECEPTORES
% ==============================================================================

% Trayectoria hacia R1.
% La flecha indica el sentido de propagación: fuente -> receptor.

\draw[pathRone]
    (NodeS) --
    node[pos=0.53, above left=2pt, font=\small]
    {$r_1$}
    (NodeR1);

% Trayectoria hacia R2.

\draw[pathRtwo]
    (NodeS) --
    node[pos=0.52, right=3pt, font=\small]
    {$r_2$}
    (NodeR2);

% ==============================================================================
% 7. DIRECCIÓN CENTRAL Y ÁNGULO DE INCIDENCIA
% ==============================================================================

% Línea de referencia desde el punto medio del arreglo hasta la fuente.

\draw[reference, ->]
    (Orig) --
    node[pos=0.63, below right=2pt, font=\small]
    {$r$}
    (S);

% Ángulo medido desde el eje normal positivo hasta la dirección central.

\draw[
    line width=0.85pt,
    draw=black!85
]
    ($(Orig)+(0,1.65)$)
    arc[start angle=90, end angle=66.2, radius=1.65];

\node[
    font=\small\bfseries
] at ($(Orig)+(78:2.0)$) {$\theta$};

% ==============================================================================
% 8. LÍMITES EXPLÍCITOS DEL DIBUJO
% ==============================================================================

% Se conserva el desplazamiento del diagrama y el límite inferior
% para alojar la cota y su etiqueta.

\path[use as bounding box]
    (-6.35, -2.15) rectangle (9.3, 7.75);

\end{tikzpicture}
